%%
%% This is file `sample-acmcp.tex',
%% generated with the docstrip utility.
%%
%% The original source files were:
%%
%% samples.dtx  (with options: `all,journal,acmcp')
%%
%% IMPORTANT NOTICE:
%%
%% For the copyright see the source file.
%%
%% Any modified versions of this file must be renamed
%% with new filenames distinct from sample-acmcp.tex.
%%
%% For distribution of the original source see the terms
%% for copying and modification in the file samples.dtx.
%%
%% This generated file may be distributed as long as the
%% original source files, as listed above, are part of the
%% same distribution. (The sources need not necessarily be
%% in the same archive or directory.)
%%
%%
%% Commands for TeXCount
%TC:macro~\cite [option:text,text]
%TC:macro~\citep [option:text,text]
%TC:macro~\citet [option:text,text]
%TC:envir table 0 1
%TC:envir table* 0 1
%TC:envir tabular [ignore] word
%TC:envir displaymath 0 word
%TC:envir math 0 word
%TC:envir comment 0 0
%%
%% The first command in your LaTeX source must be the \documentclass
%% command.
%%
%% For submission and review of your manuscript please change the
%% command to \documentclass[manuscript, screen, review]{acmart}.
%%
%% When submitting camera ready or to TAPS, please change the command
%% to \documentclass[sigconf]{acmart} or whichever template is required
%% for your publication.
%%
%%
\documentclass[sigconf, nonacm]{acmart}
\usepackage[brazilian]{babel}
\raggedbottom%
\emergencystretch=1em
\makeatletter
\@ACM@balancefalse%
\makeatother
%%
%% \BibTeX command to typeset BibTeX logo in the docs
\AtBeginDocument{%
  \providecommand\BibTeX{{%
    Bib\TeX}}}

%% Metadados de direitos/journal do template removidos: documento
%% [nonacm], fora do fluxo editorial da ACM.

%%
%% Submission ID.
%% Use this when submitting an article to a sponsored event. You'll
%% receive a unique submission ID from the organizers
%% of the event, and this ID should be used as the parameter to this command.
%%\acmSubmissionID{123-A56-BU3}




%%
%% end of the preamble, start of the body of the document source.
\begin{document}

%%
%% The "title" command has an optional parameter,
%% allowing the author to define a "short title" to be used in page headers.
\title{Reprodução Computacional da Detecção de Linguagem Tóxica em Português Brasileiro: um Estudo do ToLD-Br com Ferramentas Atuais}


%%
%% The "author" command and its associated commands are used to define
%% the authors and their affiliations.
%% Of note is the shared affiliation of the first two authors, and the
%% "authornote" and "authornotemark" commands
%% used to denote shared contribution to the research.
\author{Pedro Manoel Herminio Alves}
\email{pedro.alves@copin.ufcg.edu.br}
\affiliation{%
  \institution{Universidade Federal de Campina Grande (UFCG)}
  \city{Campina Grande}
  \state{Paraíba}
  \country{Brasil}
}



%%
%% By default, the full list of authors will be used in the page
%% headers. Often, this list is too long, and will overlap
%% other information printed in the page headers. This command allows
%% the author to define a more concise list
%% of authors' names for this purpose.
\renewcommand{\shortauthors}{P.~M.~H. Alves}
%%
%% Links to code and data
% \acmCodeLink{https://github.com/borisveytsman/acmart}
% \acmDataLink{htps://zenodo.org/link}
%%
%% Authors' contribution
% \acmContributions{BT and GKMT designed the study; LT, VB, and AP
%   conducted the experiments, BR, HC, CP and JS analyzed the results,
%   JPK developed analytical predictions, all authors participated in
%   writing the manuscript.}
%%
%% Sometimes the addresses are too long to fit on the page.  In this
%% case uncomment the lines below and fill them accodingly.
%%
%% \authorsaddresses{Corresponding author: Ben Trovato,
%% \href{mailto:trovato@corporation.com}{trovato@corporation.com};
%% Institute for Clarity in Documentation, P.O. Box 1212, Dublin,
%% Ohio, USA, 43017-6221}
%%
%%
%% Keywords. The author(s) should pick words that accurately describe
%% the work being presented. Separate the keywords with commas.
% \begin{CCSXML}
% <ccs2012>
%   <concept>
%     <concept_id>10010147.10010178.10010179</concept_id>
%     <concept_desc>Computing methodologies~Natural language processing</concept_desc>
%     <concept_significance>500</concept_significance>
%   </concept>
%   <concept>
%     <concept_id>10010147.10010257.10010258.10010259</concept_id>
%     <concept_desc>Computing methodologies~Supervised learning by classification</concept_desc>
%     <concept_significance>300</concept_significance>
%   </concept>
% </ccs2012>
% \end{CCSXML}

% \ccsdesc[500]{Computing methodologies~Natural language processing}
% \ccsdesc[300]{Computing methodologies~Supervised learning by classification}

%% --- Resumo ---
\begin{abstract}
Este trabalho reproduz computacionalmente o estudo ToLD-Br sobre detecção
de linguagem tóxica em português brasileiro, originalmente publicado com
\emph{notebooks} de 2020. O objetivo é verificar se os resultados e as
conclusões principais se mantêm em um ambiente atual, no qual parte do
\emph{stack} original --- especialmente \texttt{simple\allowbreak{}transformers},
\texttt{apex}, CUDA e \emph{auto-sklearn} --- deixou de ser diretamente
executável. Mantivemos o conjunto de dados, as tarefas e as comparações
centrais, mas modernizamos a implementação com \emph{PyTorch}, \emph{Hugging
Face Transformers} e \emph{scikit-learn}, registrando cada adaptação. Na
classificação binária, todos os modelos reproduzidos ficaram a no máximo
1,6 ponto de \emph{macro-F1} dos valores publicados, preservando as
conclusões centrais: modelos baseados em BERT superam a linha de base, o
cenário \emph{zero-shot} é substancialmente inferior e categorias raras
concentram mais \emph{false negatives}. A curva de aprendizado também
confirma a estabilização do desempenho a partir de cerca de seis mil
exemplos. A principal divergência ocorre na tarefa multirrótulo: o mBERT
reproduzido melhora em \emph{average precision}, mas com \emph{Hamming loss}
quase igual, indicando uma redistribuição entre \emph{false negatives} e
\emph{false positives}, e não uma superioridade inequívoca. Os resultados
mostram que as afirmações científicas principais do estudo são robustas,
embora a reprodução exija adaptações explícitas e rastreáveis para lidar
com a obsolescência do ambiente computacional.
\end{abstract}

\keywords{linguagem tóxica, português brasileiro, reprodutibilidade, ToLD-Br, BERT, PLN}

\maketitle

%% =====================================================================
\section{Introdução}
\label{sec:introducao}
 
As redes sociais conectam pessoas em escala global, mas também se tornaram
um canal para discurso de ódio, ofensas e ataques a grupos. Como o volume
de mensagens publicadas todos os dias é gigantesco, revisá-las manualmente
é inviável, e identificar automaticamente os comentários tóxicos tornou-se
um problema central no Processamento de Linguagem Natural (PLN) --- a área
da computação que ensina máquinas a interpretar texto. Embora esses
comentários sejam uma minoria, eles causam dano real a pessoas e
comunidades \cite{leite2020toxic}; detectá-los de forma automática é,
portanto, essencial tanto para estudar o fenômeno quanto para moderar as
plataformas.
 
A maior parte dessa pesquisa, porém, concentra-se no inglês, que dispõe de
muitos recursos --- por exemplo, o conjunto de dados OLID
\cite{zampieri2019predicting} e a série de competições \emph{OffensEval}
\cite{zampieri2019predicting}. O português
brasileiro, falado por centenas de milhões de pessoas, tem bem menos dados
rotulados e modelos próprios. Para reduzir essa lacuna,
\citet{leite2020toxic} criaram o ToLD-Br (\emph{Toxic Language Dataset for
Brazilian Portuguese}), uma coleção de 21.000 mensagens (\emph{tweets}) em
português classificadas por voluntários em seis tipos de toxicidade. Com
esse material, os autores compararam dois caminhos: um classificador
clássico e modelos modernos baseados em \emph{BERT} --- uma família de
modelos de linguagem que aprende padrões do idioma a partir de enormes
quantidades de texto e depois é especializada para uma tarefa específica
\cite{devlin2019bert}. Três conclusões se destacaram: (i) o BR-BERT obteve \emph{macro-F1} ligeiramente
superior ao M-BERT-BR, embora o M-BERT-BR tenha apresentado melhor
comportamento na classe positiva e menos \emph{false negatives}; (ii) um modelo
que só viu exemplos em inglês tem desempenho muito pior ao ser aplicado ao
português; e (iii) os tipos de toxicidade menos frequentes na base, como
racismo e xenofobia, são os mais difíceis de identificar, sendo os que
mais ``escapam'' do classificador.
 
Reexecutar esse estudo, anos depois, revela um problema que vai além de
apenas conferir números. Ao tentar reaproveitar o material publicado pelos
autores --- \emph{notebooks} de 2020, isto é, documentos de código
pensados para rodar no serviço de nuvem Google Colab ---, descobrimos que
ele não funciona mais nos computadores e sistemas atuais. O principal
obstáculo é técnico: as placas de vídeo recentes exigem versões mais novas
do software que as controla, incompatíveis com as bibliotecas em que o
código foi escrito; somam-se a isso ferramentas que foram descontinuadas e
dependências específicas do Colab. Esse é um exemplo concreto da chamada
\emph{crise de reprodutibilidade} em aprendizado de máquina \cite{pineau2021improving}: programas que
funcionavam poucos anos atrás deixam de rodar por causa da rápida evolução
do \emph{hardware} e do \emph{software}. Refazer o estudo com ferramentas
atuais serve, ao mesmo tempo, para verificar se suas conclusões continuam
válidas e para registrar o esforço necessário para mantê-lo executável.
 
Diante disso, este trabalho reproduz o desenho experimental geral
de \citet{leite2020toxic}, mantendo o ToLD-Br, as tarefas avaliadas e as
principais comparações, mas atualizando a infraestrutura e algumas decisões
de implementação necessárias para execução no ambiente atual. Em vez de
reescrever o método a partir do texto do artigo, partimos dos
\emph{notebooks} originais dos autores \cite{leite2020toldbrgithub} e os atualizamos para ferramentas
atuais --- \emph{PyTorch} \cite{paszke2019pytorch}, \emph{Hugging Face
Transformers} \cite{wolf2020transformers} e \emph{scikit-learn}
\cite{pedregosa2011scikit} ---, alterando o mínimo possível e anotando cada
mudança. Assim, esta reprodução deve ser entendida como uma reprodução
computacional com adaptações de ambiente e implementação, não como uma
repetição bit a bit do experimento original.
 
\subsection{Questões de Pesquisa}
\label{sec:rqs}
 
Três perguntas orientam o trabalho:
 
\begin{itemize}
  \item \textbf{RQ1 (os números se repetem?).} Os resultados relatados no
        artigo --- em torno de $0{,}74$ na métrica principal adotada (o
        \emph{macro-F1}, explicado na Seção~\ref{sec:avaliacao}) para o
        classificador clássico e entre $0{,}75$ e $0{,}76$ para os modelos
        baseados em \emph{BERT} treinados com dados em português ---
        reaparecem quando executamos os \emph{notebooks} atualizados sobre
        os mesmos dados?
  \item \textbf{RQ2 (as conclusões se mantêm?).} Continuam valendo as
        observações do estudo: (i) a relação observada entre BR-BERT e
        M-BERT-BR se mantém, considerando tanto o \emph{macro-F1} quanto o
        desempenho na classe positiva; (ii) treinar apenas com inglês leva
        a um desempenho bem pior; (iii) as categorias menos frequentes
        concentram mais erros; e (iv) identificar o tipo específico de
        toxicidade é mais difícil do que apenas dizer se o texto é tóxico?
  \item \textbf{RQ3 (o que mudou e quanto isso pesa?).} Quais adaptações
        foram necessárias para rodar o estudo hoje e o quanto elas afetam a
        comparação com o original?
\end{itemize}
 
%% =====================================================================
\section{Metodologia}
\label{sec:metodologia}
 
\subsection{Caracterização do Estudo}
\label{sec:caracterizacao}
 
Este é um estudo \emph{comparativo}: reexecutamos os principais
experimentos do artigo original e colocamos nossos resultados lado a lado
com os valores publicados. Como mantemos o conjunto de dados, as tarefas e
as comparações centrais, mas trocamos a infraestrutura e atualizamos parte
da implementação, trata-se de uma reprodução computacional com adaptações de
ambiente; e, como catalogamos de forma sistemática todas as mudanças feitas,
o trabalho tem também um caráter \emph{descritivo}.
 
A diferença em relação a reescrever o método do zero é importante:
partimos dos \emph{notebooks} originais dos autores (de 2020) e os
modernizamos para rodar em 2026, seguindo o princípio de \emph{alterar o
mínimo possível} e preservar ao máximo o código original. Essa
modernização não foi uma opção, e sim uma necessidade. As ferramentas de
2020 --- a biblioteca \texttt{simple\allowbreak{}transformers}, o componente
\texttt{apex} da NVIDIA e o \emph{auto-sklearn} (usado na linha de base)
--- não rodam de forma direta no ambiente atual por quatro motivos: (i) as
placas de vídeo recentes exigem uma versão de CUDA (o software que controla
a placa) igual ou superior a~11.8, enquanto o código foi feito para uma
versão antiga; (ii) o \texttt{apex} tornou-se uma dependência frágil para
reprodução em ambientes atuais, por exigir compatibilidade específica entre
PyTorch, CUDA e extensões C++/CUDA; (iii) os \emph{notebooks} dependem de
recursos específicos do Google Colab e contêm pequenos defeitos de
execução; e (iv) o \emph{auto-sklearn} só funciona diretamente em Linux.
 
O trabalho seguiu quatro etapas --- \emph{copiar}, \emph{renomear},
\emph{adaptar} e \emph{documentar} ---, pensadas para que qualquer pessoa
consiga rastrear o que mudou entre o código original e o nosso. Primeiro,
os \emph{notebooks} foram copiados para uma pasta separada,
deixando os originais intactos como referência. Em seguida, foram
renomeados com um prefixo numérico que indica a ordem de execução
(de \texttt{01\_\allowbreak generate\_\allowbreak dataset} a
\texttt{06\_\allowbreak example\_\allowbreak pretrained}); um sétimo,
\texttt{07\_\allowbreak error\_\allowbreak analysis}, foi criado por nós
para a análise de erros, que os autores não disponibilizaram em
\emph{notebook}; e um oitavo, \texttt{08\_\allowbreak compare\_\allowbreak
figures}, também criado por nós, gera as figuras que colocam lado a lado
os resultados do artigo e os desta reprodução (Figuras R.1 a R.4). Cada adaptação foi rotulada conforme o motivo ---
(T) tecnologia descontinuada, (B) correção de defeito ou
(A) fidelidade ao artigo --- e documentada em dois
lugares: uma explicação na própria célula alterada e um registro em um
catálogo central de mudanças. Nenhuma alteração foi feita ``em silêncio''.
 
Por fim, vale esclarecer o que conta como sucesso. Modelos como os
baseados em \emph{BERT} têm um componente de acaso no treinamento, e
mudanças de versão de biblioteca e de \emph{hardware} deslocam os
resultados em frações de ponto. Por isso não exigimos que os números batam
exatamente: consideramos a reprodução bem-sucedida quando recuperamos os
valores --- e, sobretudo, as conclusões --- dentro de uma margem pequena,
da ordem de 1 a 2 pontos de \emph{macro-F1}.
 
\subsection{Desenho Experimental e Variáveis}
\label{sec:desenho}
 
O experimento compara cinco abordagens, resumidas na
Tabela~\ref{tab:modelos}: uma linha de base clássica (BoW+SVM) e quatro
variações de modelos \emph{BERT}. O que muda de uma abordagem para outra
--- aquilo que controlamos no experimento --- é:
 
\begin{itemize}
  \item a abordagem em si: o classificador clássico ou um modelo
        \emph{BERT}, este treinado só em português (BERTimbau), em vários
        idiomas (mBERT), em português e inglês juntos (transferência) ou
        apenas em inglês (\emph{zero-shot});
  \item a tarefa: apenas dizer se o texto é tóxico (classificação
        binária) ou identificar quais dos seis tipos de toxicidade estão
        presentes (classificação multirrótulo);
  \item a quantidade de dados de treino, variada de 10\% a 100\%
        em um experimento à parte, a ``curva de aprendizado'';
  \item o tipo específico de toxicidade, usado na análise de
        erros.
\end{itemize}
 
Os demais ajustes de treinamento (Tabela~\ref{tab:hiperparametros}) são
mantidos fixos e iguais para os cinco modelos da tarefa binária, de modo
que as diferenças de desempenho venham das variáveis acima, e não de
configurações distintas; as duas exceções --- a curva de aprendizado e a
tarefa multirrótulo --- são indicadas na Seção~\ref{sec:protocolo}. As
métricas de classificação são resumidas principalmente pelo
\emph{macro-F1} (definido na Seção~\ref{sec:avaliacao}); examinamos também
os tipos de erro, com atenção especial aos \emph{false negatives} ---
mensagens tóxicas que o modelo deixa passar como inofensivas, o pior caso
para a moderação.
 
\begin{table}[t]
  \caption{As cinco abordagens comparadas.}
  \label{tab:modelos}
  \small
  \begin{tabular}{@{}p{0.30\linewidth}p{0.60\linewidth}@{}}
    \toprule
    \textbf{Abordagem} & \textbf{Descrição} \\
    \midrule
    BoW+SVM & Classificador clássico: conta as palavras de cada texto (``saco de palavras'') e separa tóxico de não-tóxico com um SVM. Substitui o BoW+\emph{auto-sklearn} do estudo original. \\
    BERTimbau & Modelo \emph{BERT} treinado em português, especializado no ToLD-Br. \\
    mBERT & Modelo \emph{BERT} multilíngue, especializado no ToLD-Br. \\
    mBERT (transf.) & Especializado em português e inglês juntos. \\
    mBERT (\emph{zero-shot}) & Especializado apenas em inglês. \\
    \bottomrule
  \end{tabular}
\end{table}
 
\subsection{Conjuntos de Dados}
\label{sec:dados}
 
\paragraph{ToLD-Br.} Usamos o ToLD-Br \cite{leite2020toxic}, formado por
21.000 \emph{tweets} em português, cada um avaliado por três voluntários em
seis categorias: LGBTQ+fobia, conteúdo obsceno, insulto, racismo, misoginia
e xenofobia. O repositório oficial traz o texto e os votos de cada
categoria, as avaliações individuais dos três anotadores e as divisões
oficiais dos dados. Para a tarefa de dizer apenas se o texto é tóxico,
juntamos as seis categorias em um único rótulo (tóxico ou não) pelo
critério mais inclusivo: basta um dos três avaliadores ter marcado alguma
categoria para o \emph{tweet} ser considerado tóxico --- o mesmo critério
do artigo original. Por esse critério, a base tem 9.255 mensagens tóxicas e
11.745 não-tóxicas, números que conferimos automaticamente e que coincidem
com os do artigo.
 
Adotamos as divisões oficiais do repositório: 16.800 \emph{tweets} para
treino, 2.100 para validação e 2.100 para teste. O conjunto de teste tem
1.128 mensagens não-tóxicas e 972 tóxicas --- os mesmos totais que aparecem
no artigo, o que confirmou que estávamos avaliando exatamente sobre o mesmo
conjunto. Seguindo o código original, os modelos \emph{BERT} são treinados
sobre treino e validação juntos (18.900 exemplos) e avaliados no teste
(2.100). Na tarefa multirrótulo, o arquivo com os 21.000 \emph{tweets} é
dividido em 90\% para treino e 10\% para teste, com uma semente fixa que
torna a divisão repetível.
 
\paragraph{OLID} Para os experimentos com inglês (transferência e
\emph{zero-shot}), usamos o OLID (\emph{Offensive Language Identification
Dataset}), da competição \emph{OffensEval} de 2019
\cite{zampieri2019predicting}, com \emph{tweets} em inglês marcados como
ofensivos ou não. Nós o obtivemos por meio do \emph{TweetEval}
\cite{barbieri2020tweeteval}, porque a versão atual da biblioteca de dados
não reconhece mais o nome antigo do conjunto. No código original dos
autores, os experimentos com inglês carregam o arquivo
\texttt{olid-training-v1.0.tsv}\footnote{O texto do artigo original
descreve a concatenação dos \emph{splits} de treino e teste do OLID, mas o
\emph{notebook} publicado carrega apenas o arquivo de treino; esta
reprodução segue o \emph{notebook}.}; por isso, nesta reprodução concatenamos os
splits de treino e validação do \emph{TweetEval}, totalizando 13.240
exemplos, de modo a recuperar o conjunto de treino do OLID usado pelo
\emph{notebook} original. Essa decisão privilegia a fidelidade ao artefato
computacional disponibilizado pelos autores.
 
\paragraph{Disponibilidade e limitações.} Os dados são públicos e não houve
nova coleta. O ToLD-Br é distribuído sob licença \emph{Creative Commons
BY-SA 4.0} e o OLID está disponível na plataforma \emph{Hugging Face}. Duas
limitações, porém, restringem o que pôde ser reproduzido. Primeiro, os
dados sobre o perfil dos anotadores e os identificadores dos \emph{tweets}
não são públicos (exigiriam contato com os autores); por isso, as partes do
artigo que dependem deles --- a caracterização do perfil dos voluntários e
a medida de concordância entre eles --- ficaram fora do alcance desta
reprodução. Segundo, o modelo já treinado que os autores disponibilizaram
não está mais acessível (o \emph{link} do Google Drive não permite mais o
\emph{download}); isso impediu
testar de ponta a ponta o \emph{notebook} de demonstração, cuja estrutura
foi verificada, mas não a execução.
 
\subsection{Modelos e Protocolo de Treinamento}
\label{sec:protocolo}
 
\paragraph{Linha de base.} A linha de base do estudo original combina o
``saco de palavras'' --- uma representação que descreve cada texto pela
contagem das palavras que ele contém, ignorando a ordem em que aparecem ---
com o \emph{auto-sklearn} \cite{feurer2015efficient}, uma ferramenta que
escolhe automaticamente o melhor classificador. Como o \emph{auto-sklearn}
só roda em Linux, nós o substituímos por um SVM (\emph{Support Vector
Machine}), um classificador estatístico clássico, mantendo a mesma
representação de saco de palavras. Essa troca está registrada entre os
desvios (Seção~\ref{sec:desvios}).
 
\paragraph{Modelos baseados em BERT} Tanto o modelo em português
(BERTimbau \cite{souza2020bertimbau}) quanto o multilíngue (mBERT
\cite{devlin2019bert}) passam por um \emph{ajuste fino} (\emph{fine-tuning}):
partimos de um modelo já pré-treinado em grandes volumes de texto e o
especializamos na detecção de toxicidade, mostrando-lhe os exemplos
rotulados do ToLD-Br. Os ajustes de treinamento adotados nesta reprodução
são apresentados na Tabela~\ref{tab:hiperparametros}. Quando esses valores
diferem do \emph{notebook} original, a diferença é registrada na
Seção~\ref{sec:desvios}.
 
\begin{table}[t]
  \caption{Principais ajustes de treinamento adotados nesta reprodução para os modelos baseados em \emph{BERT} na tarefa binária.}
  \label{tab:hiperparametros}
  \small
  \begin{tabular}{@{}ll@{}}
    \toprule
    \textbf{Ajuste} & \textbf{Valor} \\
    \midrule
    Passagens pelos dados (épocas) & 3 \\
    Semente aleatória              & 42 \\
    Taxa de aprendizado            & \texttt{4e-5} \\
    Aquecimento (\emph{warmup})    & 6\% dos passos \\
    Tamanho máximo do texto        & 128 \emph{tokens} \\
    Tamanho do lote (\emph{batch}) & 32 \\
    Formato numérico               & fp16 \\
    Validação durante o treino     & não \\
    \bottomrule
  \end{tabular}
\end{table}
 
\paragraph{Curva de aprendizado.} Para entender quanto os dados influenciam
o resultado, treinamos o modelo multilíngue com quantidades crescentes de
exemplos, de 10\% a 100\% do conjunto de treino. Cada fração é treinada
três vezes (com sementes diferentes) e reportamos a média, para reduzir o
efeito do acaso. Seguindo o \emph{notebook} original, cada uma dessas
repetições treina por uma única época --- e não as três da
Tabela~\ref{tab:hiperparametros} ---, o que torna viáveis as trinta
execuções ao custo de maior instabilidade entre repetições (visível na
Seção~\ref{sec:res-curva}).
 
\paragraph{Classificação multirrótulo.} Na tarefa de identificar o tipo de
toxicidade, o modelo multilíngue produz seis respostas independentes (uma
por categoria); consideramos uma categoria presente quando a confiança do
modelo ultrapassa 50\%. Nesta tarefa, o lote de treino é de 16 exemplos
(o código original usava 8; os demais ajustes seguem a
Tabela~\ref{tab:hiperparametros}). A linha de base correspondente treina um
classificador separado para cada categoria.
 
\subsection{Métricas e Procedimento de Avaliação}
\label{sec:avaliacao}
 
Avaliamos os modelos com o \emph{scikit-learn}~\cite{pedregosa2011scikit}.
Para a tarefa binária, usamos os nomes clássicos das métricas de
classificação: \emph{precision}, \emph{recall}, \emph{F1-score} e
\emph{macro-F1}. A \emph{precision} mede, entre as mensagens previstas
como tóxicas, quantas realmente eram tóxicas; o \emph{recall} mede, entre
as mensagens tóxicas, quantas foram encontradas pelo modelo; e o
\emph{F1-score} combina \emph{precision} e \emph{recall} em um único
valor. A métrica principal é o \emph{macro-F1}, isto é, a média não
ponderada do \emph{F1-score} das classes tóxica e não-tóxica.
Complementamos a análise com a \emph{confusion matrix} e com a
\emph{false negative rate} (FNR) por categoria. Na tarefa multirrótulo,
reportamos ainda \emph{Hamming loss} e \emph{average precision}, que
resumem, respectivamente, a fração de rótulos errados e a qualidade geral
das previsões.
 
\subsection{Ferramentas e Ambiente de Execução}
\label{sec:ambiente}
 
Todo o trabalho foi executado em um computador local --- e não na nuvem,
como no estudo original. O ambiente é o Windows~11, com uma placa de vídeo
NVIDIA GeForce RTX~4070~Ti~SUPER (16~GB) e CUDA~12.4. É justamente essa
placa recente que motiva a modernização: sua arquitetura exige uma versão
de CUDA incompatível com as bibliotecas de 2020 (Seção~\ref{sec:caracterizacao}).
 
A linguagem é o Python. Os modelos \emph{BERT} são ajustados com
\emph{PyTorch} \cite{paszke2019pytorch} e a biblioteca \emph{Hugging Face
Transformers} \cite{wolf2020transformers}; a linha de base e o cálculo das
medidas usam o \emph{scikit-learn} \cite{pedregosa2011scikit}; e a
manipulação de dados e os gráficos ficam por conta de bibliotecas usuais do
ecossistema Python, como \emph{pandas}, \emph{numpy} e \emph{matplotlib}.
Os \emph{notebooks} são executados em um ambiente virtual isolado; as
versões efetivamente utilizadas de todas as bibliotecas estão congeladas
em um arquivo de \emph{lock} (\texttt{requirements\allowbreak-lock.txt})
que acompanha o material de reprodução, para que o próprio experimento
seja fácil de repetir.
 
\subsection{Adaptações e Desvios em Relação ao Estudo Original}
\label{sec:desvios}
 
Reproduzir o estudo em um ambiente diferente exigiu algumas adaptações.
Classificamos cada uma pelo motivo --- (T) tecnologia
descontinuada, (B) correção de defeito ou (A) fidelidade
ao artigo --- e todas estão registradas no catálogo de mudanças. As
principais são:
 
\begin{enumerate}
  \item \textbf{Linha de base \textnormal{(T/A)}.} Como o
        \emph{auto-sklearn} só roda em Linux, trocamos a linha de base por
        um SVM, mantendo a mesma representação de saco de palavras.
  \item \textbf{Ferramentas de treinamento \textnormal{(T)}.} A combinação
        \texttt{simple\allowbreak{}transformers} + \texttt{apex}, que se
        mostrou incompatível com o ambiente atual de execução, foi
        substituída por uma implementação baseada na biblioteca
        \emph{Hugging Face Transformers}~\cite{wolf2020transformers}.
  \item \textbf{Fidelidade dos modelos ao artigo \textnormal{(A)}.} O
        código publicado usa um modelo mais leve, mas o texto do artigo
        descreve o BERTimbau \cite{souza2020bertimbau} e o mBERT
        \cite{devlin2019bert}; seguimos o artigo. Esta mudança altera a
        implementação dos modelos BERT em relação ao \emph{notebook}
        publicado. Além dela, a linha de base também foi alterada de
        BoW+\emph{auto-sklearn} para BoW+SVM, e os hiperparâmetros
        operacionais foram adaptados ao ambiente atual.
  \item \textbf{Aquecimento da taxa de aprendizado \textnormal{(B/T)}.}
        Reintroduzimos o ``aquecimento'' (\emph{warmup}) que era padrão na
        ferramenta original; sem ele, o modelo de transferência falhava,
        classificando quase tudo como a categoria mais comum.
  \item \textbf{Tamanho do texto e do lote \textnormal{(T)}.} Reduzimos o
        tamanho máximo do texto de 512 para 128 unidades (\emph{tokens})
        --- suficiente para um \emph{tweet} e cerca de quatro vezes mais
        rápido --- e o lote de treino de 50 para 32, por limite de memória
        da placa. Na tarefa multirrótulo, o lote passou de 8 (código
        original) para 16.
  \item \textbf{Obtenção do OLID \textnormal{(T/A)}.}
        Obtivemos o OLID por meio do \emph{TweetEval}
        \cite{barbieri2020tweeteval}, já que a biblioteca de dados atual
        não reconhece mais o nome antigo do conjunto. Para preservar a
        composição usada no \emph{notebook} original, concatenamos os
        splits de treino e validação do \emph{TweetEval}, recuperando os
        13.240 exemplos correspondentes ao arquivo
        \texttt{olid-training-v1.0.tsv}. Assim, essa adaptação diz respeito
        à forma de acesso ao dado, e não à composição efetiva do conjunto
        usado nos experimentos.
  \item \textbf{Correção de pequenos defeitos \textnormal{(B)}.} Corrigimos
        quatro defeitos de código nos \emph{notebooks}: três impediam a
        execução (duas variáveis não definidas e um nome de argumento
        trocado) e um era um erro de indentação silencioso no
        pré-processamento, fora do caminho usado pelo artigo. O
        \emph{notebook} original da linha de base continha ainda um erro
        de digitação bloqueante (chamada a um método inexistente), cuja
        correção ficou absorvida pela troca descrita no item 1. Na análise
        de erros, padronizamos o texto (maiúsculas e minúsculas, acentos e
        espaços) para casar corretamente as mensagens entre os arquivos.
\end{enumerate}
 
%% =====================================================================

\section{Resultados e Discussão}
\label{sec:resultados}
 
Esta seção apresenta os resultados em quatro blocos: classificação
binária, análise de erro por categoria, curva de aprendizado e
classificação multirrótulo. Em todos eles, a leitura combina dois níveis:
primeiro, a proximidade numérica em relação ao artigo original; segundo, a
manutenção (ou não) das conclusões qualitativas. Como os modelos neurais
podem variar levemente entre execuções e versões de biblioteca, consideramos
bem-sucedida uma reprodução que permaneça dentro de uma margem prática de
ruído --- cerca de 1 a 2 pontos de \emph{macro-F1}. No ambiente usado aqui, a
maior parte do \emph{pipeline} mostrou-se determinística: ao fixar a semente
aleatória (valor inicial que controla operações pseudoaleatórias, como
embaralhamento dos dados e inicialização do treino) e congelar as versões das
bibliotecas, uma reexecução completa recuperou dígito a dígito os cinco modelos
binários e os dez pontos da curva de aprendizado. A exceção é a tarefa
multirrótulo: nela, duas execuções do mesmo código, com a mesma semente e na
mesma máquina, produziram resultados diferentes (\emph{average precision} de
0,276 e 0,305). Essa instabilidade é coerente com o desbalanceamento extremo das
classes raras discutido na Seção~\ref{sec:res-achados}, e os números
multirrótulo aqui reportados devem ser lidos como uma amostra, não como um valor
exatamente reproduzível.
 
\subsection{Classificação binária}
\label{sec:res-binario}
 
A Tabela~\ref{tab:macrof1} e a Figura~\ref{fig:r1} mostram o resultado
principal da reprodução: todos os modelos ficam dentro da margem esperada em
relação ao artigo. O maior desvio absoluto é de 1,6 ponto de \emph{macro-F1}
(M-BERT-BR e \emph{zero-shot}); nos demais casos, a diferença fica abaixo de
1 ponto. Mais importante que a coincidência exata é a ordem dos resultados:
os modelos treinados com dados em português permanecem na faixa de 0,75 a
0,77, o \emph{baseline} fica próximo de 0,73 e o cenário \emph{zero-shot}
continua claramente inferior, com 0,544. Assim, a conclusão central do
artigo para a tarefa binária é preservada.
 
\begin{table}[t]
  \caption{Macro-F1 na classificação binária: artigo \emph{versus} reprodução.}
  \label{tab:macrof1}
  \small
  \begin{tabular}{@{}lccc@{}}
    \toprule
    \textbf{Modelo} & \textbf{Artigo} & \textbf{Reprod.} & \textbf{$\Delta$} \\
    \midrule
    \emph{Baseline} BoW (AutoML/SVM) & 0,74 & 0,730 & $-0{,}010$ \\
    BR-BERT (BERTimbau)              & 0,76 & 0,765 & $+0{,}005$ \\
    M-BERT-BR (mBERT)                & 0,75 & 0,766 & $+0{,}016$ \\
    M-BERT (transf.)                 & 0,76 & 0,752 & $-0{,}008$ \\
    M-BERT (\emph{zero-shot})        & 0,56 & 0,544 & $-0{,}016$ \\
    \bottomrule
  \end{tabular}
\end{table}
 
\begin{figure}[t]
  \centering
  \includegraphics[width=\columnwidth]{comparison/compare_macro_f1_bars.png}
  \caption{Macro-F1 na classificação binária: artigo (cinza) \emph{versus} reprodução (azul). (Figura R.1.)}
  \Description{Gráfico de barras que compara o macro-F1 reportado no artigo e o macro-F1 obtido na reprodução para os cinco modelos da tarefa binária.}
  \label{fig:r1}
\end{figure}
 
O detalhamento por classe (Tabela~\ref{tab:porclasse}) ajuda a interpretar
o empate entre BR-BERT e M-BERT-BR. Os dois têm \emph{macro-F1} praticamente
igual (0,765 e 0,766), mas erram de formas diferentes. O M-BERT-BR tem maior
\emph{recall} para a classe tóxica (0,795 contra 0,746), ou seja, deixa
passar menos mensagens tóxicas como não-tóxicas. Essa diferença explica por
que o artigo trata o M-BERT-BR como melhor na análise de erro, mesmo sem uma
vantagem expressiva no \emph{macro-F1}. A Figura~\ref{fig:r2} mostra as
\emph{confusion matrices} completas, lado a lado com as originais.

\begin{figure}[t]
  \centering
  \includegraphics[width=\columnwidth,height=0.85\textheight,keepaspectratio]{comparison/compare_fig2_binary_cm.png}
  \caption{\emph{Confusion matrices} na classificação binária: artigo (à esquerda) \emph{versus} reprodução (à direita). (Figura R.2.)}
  \Description{Conjunto de confusion matrices que compara, para cada modelo binário, as predições reportadas no artigo com as obtidas na reprodução.}
  \label{fig:r2}
\end{figure}

\begin{table*}[t]
  \caption{\emph{Precision} (P), \emph{recall} (R) e \emph{F1-score} por classe na reprodução (classe 0 = não-tóxico; classe 1 = tóxico).}
  \label{tab:porclasse}
  \small
  \begin{tabular}{@{}lcccccc@{}}
    \toprule
     & \multicolumn{3}{c}{\textbf{Classe 0 (não-tóxico)}} & \multicolumn{3}{c}{\textbf{Classe 1 (tóxico)}} \\
    \cmidrule(lr){2-4}\cmidrule(lr){5-7}
    \textbf{Modelo} & P & R & F1 & P & R & F1 \\
    \midrule
    BoW+SVM                   & 0,725 & 0,817 & 0,769 & 0,752 & 0,641 & 0,692 \\
    BR-BERT                   & 0,782 & 0,784 & 0,783 & 0,748 & 0,746 & 0,747 \\
    M-BERT-BR                 & 0,808 & 0,742 & 0,774 & 0,727 & 0,795 & 0,759 \\
    M-BERT (transf.)          & 0,795 & 0,726 & 0,759 & 0,711 & 0,783 & 0,745 \\
    M-BERT (\emph{zero-shot}) & 0,581 & 0,830 & 0,683 & 0,607 & 0,305 & 0,405 \\
    \bottomrule
  \end{tabular}
\end{table*}
 
\subsection{Análise de erro por categoria}
\label{sec:res-erro}
 
A Tabela~\ref{tab:fnrate} avalia onde o classificador falha com mais
frequência. Para cada categoria, reportamos a \emph{false negative rate} (FNR),
isto é, a fração de mensagens daquela categoria que o modelo classificou
como não-tóxicas. O padrão do artigo aparece novamente: categorias raras
continuam sendo as mais difíceis. Racismo e xenofobia têm as maiores taxas
(0,41 e 0,47), enquanto categorias mais frequentes, como obsceno e insulto,
ficam em 0,17.

Os valores não coincidem exatamente com os do artigo por uma razão
operacional. As categorias finas não estão no mesmo arquivo das predições do
conjunto de teste; por isso, foi necessário ligar os registros pelo texto da
mensagem. Para aumentar a chance de correspondência, normalizamos os textos
(minúsculas, remoção de acentos e espaços repetidos). Esse procedimento
recupera as 2.100 mensagens do teste, mas pode fundir mensagens originalmente
diferentes e ainda deixa cerca de dez linhas com codificação corrompida.
Assim, os denominadores por categoria ficam muito próximos dos do artigo,
mas não perfeitamente idênticos (ver Seção~\ref{sec:ameacas}).

 
\begin{table}[t]
  \caption{\emph{False negative rate} (FNR) por categoria: artigo \emph{versus} reprodução.}
  \label{tab:fnrate}
  \small
  \begin{tabular}{@{}lcc@{}}
    \toprule
    \textbf{Categoria} & \textbf{Artigo} & \textbf{Reprod.} \\
    \midrule
    LGBTQ+fobia & 0,20 & 0,20 \\
    Insulto     & 0,15 & 0,17 \\
    Xenofobia   & 0,68 & 0,47 \\
    Misoginia   & 0,15 & 0,16 \\
    Obsceno     & 0,17 & 0,17 \\
    Racismo     & 0,47 & 0,41 \\
    \bottomrule
  \end{tabular}
\end{table}
 
\subsection{Curva de aprendizado}
\label{sec:res-curva}
 
A curva de aprendizado responde à pergunta de quanto dado é necessário para
obter desempenho estável. A Tabela~\ref{tab:curva} e a
Figura~\ref{fig:r3b} mostram que o M-BERT-BR melhora rapidamente nas
primeiras frações do treino e passa a variar menos a partir de 30--40\% dos
dados (aproximadamente 5,7 a 7,6 mil exemplos). Isso reproduz a conclusão do
artigo de que alguns milhares de exemplos já são suficientes para chegar à
faixa de desempenho observada no conjunto completo.

O ganho de dados também muda o tipo de erro. O \emph{recall} da classe
tóxica cresce de 0,41 para 0,82, indicando que o modelo passa a deixar menos
comentários tóxicos escaparem. Em contrapartida, o \emph{recall} da classe
não-tóxica cai de 0,80 para 0,69, sinal de que o modelo se torna menos
conservador. A queda isolada em 60\% não altera a tendência geral: ela vem
de uma das três repetições, cuja perda de treino ficou estagnada enquanto as
outras duas convergiram. Trata-se, portanto, de instabilidade pontual do
\emph{fine-tuning} com apenas uma época, não de uma mudança de conclusão.
 
\begin{table}[t]
  \caption{Curva de aprendizado (M-BERT-BR; média de três repetições por fração).}
  \label{tab:curva}
  \small
  \begin{tabular}{@{}cccc@{}}
    \toprule
    \textbf{\% treino} & \textbf{macro-F1} & \textbf{recall tóxica} & \textbf{recall não-tóx.} \\
    \midrule
    10\%  & 0,584 & 0,409 & 0,798 \\
    20\%  & 0,679 & 0,689 & 0,673 \\
    30\%  & 0,708 & 0,698 & 0,721 \\
    40\%  & 0,720 & 0,747 & 0,698 \\
    50\%  & 0,734 & 0,776 & 0,699 \\
    60\%  & 0,645 & 0,577 & 0,769 \\
    70\%  & 0,740 & 0,791 & 0,695 \\
    80\%  & 0,736 & 0,806 & 0,676 \\
    90\%  & 0,747 & 0,799 & 0,702 \\
    100\% & 0,749 & 0,818 & 0,689 \\
    \bottomrule
  \end{tabular}
\end{table}
 
\begin{figure*}[t]
  \centering
  \includegraphics[width=\textwidth]{comparison/compare_fig3b_overlay_learning_curve.png}
  \caption{Curva de aprendizado com as curvas reais sobrepostas, eixo em \% dos dados de treino (linha sólida = reprodução; tracejada = artigo). As curvas têm a mesma forma e convergem a 100\%. A única diferença de base está no ponto de 100\%: no artigo, ele corresponde ao conjunto completo (21.000 exemplos, incluindo o teste); na reprodução, ao conjunto de treino (18.900, com o teste sempre separado). (Figura R.3b.)}
  \Description{Curva de aprendizado que mostra a evolução do macro-F1 conforme aumenta a fração de dados de treino, comparando a reprodução com a curva do artigo.}
  \label{fig:r3b}
\end{figure*}
 
Há uma diferença de protocolo apenas no último ponto da curva. No artigo, o
ponto de 100\% usa os 21.000 exemplos do ToLD-Br, incluindo o teste; nesta
reprodução, todos os pontos, inclusive 100\%, usam apenas os 18.900 exemplos
de treino e validação, mantendo o teste separado. Essa escolha evita
vazamento de dados. Como as curvas sobrepostas na Figura~\ref{fig:r3b} têm a
mesma forma e convergem no ponto final, essa diferença de escala não afeta a
interpretação do resultado.
 
\subsection{Classificação multirrótulo}
\label{sec:res-multi}
 
A tarefa multirrótulo é o único bloco em que os números não acompanham
integralmente o artigo. A linha de base permanece muito próxima
(\emph{average precision} 0,209 contra 0,20; \emph{Hamming loss} 0,069
contra 0,08), mas o mBERT muda de comportamento: sua \emph{average
precision} sobe de 0,19 no artigo para 0,276 nesta reprodução. Ao mesmo
tempo, a \emph{Hamming loss} fica praticamente igual (0,067 contra 0,07), o
que indica que o volume total de erros não mudou muito. A diferença está em
como esses erros se distribuem entre \emph{false negatives} e \emph{false positives},
como discutido na Seção~\ref{sec:res-achados}. A Figura~\ref{fig:r4} mostra
as \emph{confusion matrices} por rótulo.

\begin{figure}[t]
  \centering
  \includegraphics[width=\columnwidth,height=0.85\textheight,keepaspectratio]{comparison/compare_fig4_multilabel_cm.png}
  \caption{\emph{Confusion matrices} por rótulo na tarefa multirrótulo (mBERT): artigo (à esquerda) \emph{versus} reprodução (à direita). (Figura R.4.)}
  \Description{Confusion matrices por categoria de toxicidade na tarefa multirrótulo, comparando os resultados do artigo com os resultados reproduzidos.}
  \label{fig:r4}
\end{figure}

\begin{table}[t]
  \caption{Classificação multirrótulo: artigo \emph{versus} reprodução.}
  \label{tab:multi}
  \small
  \begin{tabular}{@{}llccc@{}}
    \toprule
    \textbf{Modelo} & \textbf{Métrica} & \textbf{Artigo} & \textbf{Reprod.} & \textbf{$\Delta$} \\
    \midrule
    BoW+SVM & \emph{Hamming loss} & 0,08 & 0,069 & $-0{,}011$ \\
    BoW+SVM & \emph{Average precision} & 0,20 & 0,209 & $+0{,}009$ \\
    mBERT   & \emph{Hamming loss} & 0,07 & 0,067 & $-0{,}003$ \\
    mBERT   & \emph{Average precision} & 0,19 & \textbf{0,276} & $\mathbf{+0{,}086}$ \\
    \bottomrule
  \end{tabular}
\end{table}
 
\subsection{Relação com as hipóteses}
\label{sec:res-hipoteses}
 
Para ligar os resultados às questões de pesquisa, a
Tabela~\ref{tab:hipoteses} resume oito afirmações derivadas do artigo
original. Consideramos uma afirmação confirmada quando a tendência descrita
no artigo se mantém nos nossos números, mesmo que os valores não sejam
idênticos. O balanço geral é positivo: seis afirmações são confirmadas, uma
fica parcial (o suposto ganho do modelo monolíngue vira empate dentro do
ruído) e uma é rejeitada (a vantagem do \emph{baseline} sobre BERT na
\emph{average precision} multirrótulo).
 
\begin{table*}[t]
  \caption{Afirmações do artigo testadas na reprodução e respectivo veredito.}
  \label{tab:hipoteses}
  \small
  \begin{tabular}{@{}cp{0.40\linewidth}lp{0.30\linewidth}@{}}
    \toprule
     & \textbf{Afirmação do artigo} & \textbf{Veredito} & \textbf{Evidência (reproduzida)} \\
    \midrule
    H1 & BERT supera o \emph{baseline} na tarefa binária & Confirmada & BR-BERT 0,765 / M-BERT-BR 0,766 / transf.\ 0,752 $>$ 0,730 \\
    H2 & Modelo monolíngue um pouco melhor que o multilíngue & Parcial (empate) & BR-BERT 0,765 $\approx$ M-BERT-BR 0,766 (dentro do ruído) \\
    H3 & Transferência não melhora sobre o monolíngue & Confirmada & transf.\ 0,752 $<$ 0,766 e 0,765 \\
    H4 & \emph{Zero-shot} é bem inferior & Confirmada & 0,544; \emph{recall} da classe tóxica 0,305 \\
    H5 & Mais \emph{false negatives} nas classes menores & Confirmada & racismo 0,41; xenofobia 0,47; obsceno/insulto 0,17 \\
    H6 & Pelo menos cerca de 6 mil exemplos são necessários & Confirmada & estabiliza em 30--40\% (cerca de 5,7--7,6 mil) \\
    H7 & Multirrótulo difícil; classes raras quase tudo negativo & Confirmada & \emph{true positives}: racismo 0; xenofobia 1; misoginia 11 \\
    H8 & \emph{Baseline} melhor que BERT no multirrótulo (\emph{average precision}) & Rejeitada & mBERT 0,276 $>$ \emph{baseline} 0,209 \\
    \bottomrule
  \end{tabular}
\end{table*}
 
\subsection{Resultado inesperado e principais achados}
\label{sec:res-achados}
 
A principal divergência está no multirrótulo. À primeira vista, o mBERT
reproduzido parece melhor que o do artigo, pois a \emph{average precision}
sobe de 0,19 para 0,276. Essa leitura, porém, é incompleta. A \emph{Hamming
loss} quase não muda (0,067 contra 0,07), indicando que a quantidade total
de erros é semelhante. O que muda é o tipo de erro: nas categorias mais
frequentes, o modelo reproduzido prevê mais positivos (por exemplo, 484
\emph{true positives} em obsceno, contra 205 no artigo), mas também gera
mais \emph{false positives} (291 em obsceno, contra 38). Portanto, o resultado
não indica superioridade geral; indica um \emph{trade-off} entre maior
\emph{recall} e menor \emph{precision}.

A explicação mais plausível é a mudança no mecanismo de treino. O modelo e
os dados são os mesmos descritos no artigo, e a linha de base praticamente
coincide. A diferença principal está no \emph{stack} de \emph{mixed precision}: o
código de 2020 dependia do \texttt{apex}, enquanto a reprodução usa o
mecanismo moderno do PyTorch. Registros do treino original indicam estouros
de gradiente e passos do otimizador ignorados no início, o que tende a
produzir um modelo mais conservador. Com o treino moderno, o modelo aprende
mais exemplos positivos; isso aumenta a \emph{average precision}, mas também
eleva os \emph{false positives}. O efeito aparece no multirrótulo porque algumas
classes são extremamente raras (racismo tem cerca de 11 positivos em 2.100
exemplos), tornando muito atraente para o modelo prever quase tudo como
negativo. Na tarefa binária, em que cerca de 44\% do teste é positivo, esse
problema é muito menor.

Em síntese, a reprodução sustenta as conclusões principais do artigo:
\begin{enumerate}
  \item a classificação binária, resultado central do estudo, foi reproduzida
        com sucesso, com todos os modelos a no máximo 1,6 ponto de
        \emph{macro-F1};
  \item modelos baseados em BERT continuam superando o \emph{baseline}, e o
        cenário \emph{zero-shot} continua substancialmente inferior;
  \item categorias raras seguem concentrando \emph{false negatives}, e a curva de
        aprendizado confirma que cerca de 6 mil exemplos já levam a
        resultados estáveis;
  \item a vantagem do modelo monolíngue sobre o multilíngue não se confirmou,
        aparecendo como empate dentro da margem de ruído;
  \item a única divergência substantiva é o multirrótulo, explicado como
        redistribuição entre \emph{false negatives} e \emph{false positives}, não como
        ganho inequívoco do modelo reproduzido.
\end{enumerate}
 
%% =====================================================================
\section{Ameaças à Validade}
\label{sec:ameacas}
 
\paragraph{Validade de construto e interna.} Três escolhas afetam a
comparação direta com o artigo. A linha de base é diferente: usamos BoW+SVM
no lugar do BoW+\emph{auto-sklearn} (que só executa em Linux); ambas as opções
ficam em torno de 0,73 de \emph{macro-F1}, mas não são o mesmo
classificador. Além disso, alguns hiperparâmetros não são totalmente
especificados no artigo (taxa de aprendizado, \emph{warmup} e tamanho do
lote); adotamos os valores padrão da ferramenta usada no código original, o
que é uma inferência. Por fim, a análise de erro depende de uma normalização
de texto para alinhar as predições às categorias finas, e cerca de dez
linhas têm codificação corrompida irrecuperável; por isso, os totais por
categoria ficam próximos, mas não idênticos aos do artigo (por exemplo,
insulto com 444 exemplos contra 448).
 
\paragraph{Validade estatística e de conclusão.} Exceto na curva de
aprendizado (três repetições por ponto), os modelos binários foram
treinados com uma única semente, e não computamos testes de significância
formais. Assim, ao falar em ``diferença'' entre métodos, referimo-nos à
magnitude prática diante da margem de ruído de 1 a 2 pontos, e não a
significância estatística. A divergência do multirrótulo é explicada de
forma qualitativa (\emph{mixed precision} moderna \emph{versus} o \texttt{apex} de
2020); confirmar a causa exigiria reexecutar exatamente o \emph{stack} de
2020, o que está fora do escopo.
 
\paragraph{Validade externa e reprodutibilidade.} As versões de bibliotecas
e de CUDA são muito mais recentes que as de 2020; em outras máquinas, os
resultados dos modelos \emph{BERT} podem variar em 1 a 2 pontos. Neste ambiente,
com semente fixa, uma reexecução completa reproduziu dígito a dígito a tarefa
binária e a curva de aprendizado, mas \emph{não} a tarefa multirrótulo, cuja
\emph{average precision} variou entre execuções (0,276 e 0,305) sem nenhuma
mudança de código, semente ou biblioteca --- provável efeito de operações
não-determinísticas de GPU, amplificado pelas classes raras
(Seção~\ref{sec:res-achados}). Portanto, as conclusões multirrótulo devem
apoiar-se na direção do resultado, e não no valor exato. O
ambiente é Windows com uma única GPU, e o \emph{baseline} AutoML do artigo
só executa em Linux. O OLID foi obtido pelo \emph{TweetEval}, cujo texto passa
por uma leve normalização (menções viram \texttt{@user}, \emph{links} viram
\texttt{http}); é equivalente, mas não idêntico byte a byte ao arquivo
original.
 
\paragraph{Cobertura.} Dois itens do artigo ficaram fora de escopo. A
concordância entre anotadores (coeficiente de Krippendorff) não foi
reproduzida numericamente, pois o cálculo é feito por uma ferramenta
externa ao \emph{pipeline} de \emph{notebooks}. E a análise de robustez do
multirrótulo (variando semente e hiperparâmetros), que ajudaria a fechar a
causa da divergência, não faz parte do \emph{pipeline} e não foi
preservada.
 
%% =====================================================================
\section{Conclusões}
\label{sec:conclusoes}
 
Esta reprodução confirma o resultado central do estudo de
\citet{leite2020toxic}. Na classificação binária --- a tarefa principal ---,
os cinco modelos reproduziram o artigo dentro da margem de ruído (no máximo
1,6 ponto de \emph{macro-F1}), e as conclusões qualitativas centrais se
sustentam: os modelos baseados em \emph{BERT} superam o \emph{baseline}
clássico, a transferência a partir do inglês não melhora sobre o modelo
monolíngue, o cenário \emph{zero-shot} é muito inferior, as categorias
minoritárias concentram os \emph{false negatives} e o desempenho estabiliza a
partir de cerca de seis mil exemplos.
 
Foi possível modernizar integralmente o ambiente --- do Google Colab com
\texttt{simple\allowbreak{}transformers} e \texttt{apex}, de 2020, para o
\emph{Hugging Face Transformers} em 2026 --- preservando os modelos e o
protocolo do artigo, com alterações mínimas e rastreáveis sobre os
\emph{notebooks} dos autores. Isso sugere que reproduções a partir dos
artefatos originais, e não apenas do texto, são viáveis mesmo após a
obsolescência do \emph{software}, desde que cada adaptação seja documentada.
 
A única divergência está no multirrótulo: o mBERT reproduzido supera o do
artigo em \emph{average precision} (0,276 contra 0,19), mas como um
\emph{trade-off} --- mais \emph{recall}, mais \emph{false positives} e \emph{Hamming
loss} praticamente igual ---, atribuível à maior estabilidade do mecanismo
moderno de \emph{mixed precision} frente ao \texttt{apex} de 2020, efeito que só se
manifesta sob o desbalanceamento extremo das classes raras. A vantagem do
modelo monolíngue sobre o multilíngue não se confirmou, ficando num empate
dentro do ruído, o que é coerente com o próprio artigo, que elege o modelo
multilíngue treinado em português como o melhor na análise de erro.
 
No conjunto, trata-se de uma reprodução bem-sucedida das afirmações
científicas principais do artigo, com uma divergência compreendida e
documentada e com as limitações honestamente registradas. Como trabalho
futuro, fechar a causa da divergência do multirrótulo exigiria reexecutar o
\emph{stack} numérico de 2020 e conduzir uma análise de robustez por
semente.

%% =====================================================================
\section{Reprodutibilidade (artefatos)}
\label{sec:reprodutibilidade}

\paragraph{Onde estão os artefatos.} Todo o material desta reprodução --- dados,
\emph{notebooks} modernizados, catálogo de mudanças e relatório --- está reunido em um
repositório público no GitHub:
\begin{center}
\url{https://github.com/Pedro-Manoel/ToLD-Br-Reproducao}
\end{center}
O repositório inclui: o conjunto ToLD-Br (\texttt{ToLD-BR.csv} e
\texttt{ToLD-BR\_alpha.csv}, sob licença \emph{Creative Commons BY-SA 4.0}); os
\emph{notebooks} da reprodução em \texttt{reproduction/notebooks/}, numerados pela ordem de
execução; o catálogo de todas as adaptações em \texttt{reproduction/notebooks/CHANGES.md}
(rotuladas T/B/A, conforme a Seção~\ref{sec:desvios}); o relatório detalhado
\texttt{reproduction/REPORT.md}; e o ambiente congelado em
\texttt{reproduction/requirements\allowbreak-lock.txt}.

\paragraph{Como reproduzir.} A reprodução roda inteiramente pelos \emph{notebooks}. Após
clonar o repositório, prepara-se um ambiente virtual isolado --- em Windows com
\texttt{setup\_env.ps1}, em Linux/macOS com \texttt{setup\_env.sh}, ou de forma
multiplataforma com \emph{conda} a partir de \texttt{environment.yml} --- e registra-se o
\emph{kernel} Jupyter \texttt{toldbr-repro}. Os \emph{notebooks} do \emph{pipeline}
principal são então executados na ordem
\texttt{01}$\rightarrow$\texttt{02}$\rightarrow$\texttt{03}$\rightarrow$\texttt{07}$\rightarrow$\texttt{08}:
o \texttt{03} produz os cinco modelos binários, a curva de aprendizado e a tarefa
multirrótulo; o \texttt{07}, a análise de erro por categoria; e o \texttt{08}, as figuras
comparativas deste artigo. Os \emph{notebooks} \texttt{04} (curva), \texttt{05} (nuvens de
palavras) e \texttt{06} (demonstração do modelo pré-treinado, cujo \emph{link} original está
fora do ar) são complementares e não geram os resultados aqui reportados. O tempo total é de
aproximadamente 1,5 a 2 horas, dominado pelo \texttt{03}. As instruções completas, para os
dois sistemas operacionais, estão no \texttt{README} do repositório.

\paragraph{Recursos necessários.} O ambiente de referência usa Python~3.13 e uma GPU NVIDIA
com CUDA (RTX~4070~Ti~SUPER, 16~GB); as bibliotecas principais são \emph{PyTorch}, \emph{Hugging
Face Transformers} e \emph{scikit-learn} (versões exatas em
\texttt{requirements\allowbreak-lock.txt}). O \emph{pipeline} também executa em CPU --- os
\emph{notebooks} ativam \emph{mixed precision} apenas quando há GPU disponível ---, embora o
treino dos modelos \emph{BERT} (\texttt{03}) fique consideravelmente mais lento. O ToLD-Br
acompanha o repositório; o OLID é obtido em tempo de execução pelo \emph{TweetEval}
(\emph{Hugging Face}), sem \emph{download} manual. Fixadas a semente aleatória e as versões das
bibliotecas, uma reexecução completa no ambiente de referência recupera dígito a dígito a tarefa
binária e a curva de aprendizado; a tarefa multirrótulo, porém, varia entre execuções mesmo com
a mesma semente (Seções~\ref{sec:res-achados} e~\ref{sec:ameacas}). Em outro \emph{hardware},
variações da ordem de 1 a 2 pontos de \emph{macro-F1} são esperadas.

%% =====================================================================
\section{Declaração de Uso de Inteligência Artificial}
\label{sec:declaracao-ia}

Em atenção às boas práticas de integridade e transparência na pesquisa,
declara-se que na elaboração deste trabalho foi utilizada a ferramenta de
inteligência artificial generativa \emph{Claude}, da Anthropic, no modelo
\emph{Opus~4.8} (identificador \texttt{claude\allowbreak-opus\allowbreak-4\allowbreak-8}).

O uso deu-se em duas frentes, sempre como auxílio: (i) na \textbf{revisão e
validação do código} da reprodução --- conferência dos \emph{notebooks}
modernizados, verificação da consistência entre os artefatos gerados e os
valores relatados, e apoio na execução e na checagem do \emph{pipeline}; e (ii)
na \textbf{revisão e validação da escrita} --- revisão de clareza, coerência e
consistência interna do texto.

O emprego da ferramenta restringiu-se a esse apoio instrumental. A concepção do
estudo, as decisões metodológicas, a condução dos experimentos, a interpretação
dos resultados e as conclusões são de responsabilidade exclusiva do autor, que
revisou e validou integralmente todo o conteúdo aqui apresentado. A ferramenta
não é autora do trabalho, não lhe é atribuída autoria ou coautoria, e não
respondeu por nenhuma decisão científica. Nenhum resultado numérico foi gerado
ou estimado pela ferramenta: todos os valores reportados provêm da execução dos
\emph{notebooks} descritos na Seção~\ref{sec:reprodutibilidade}.

%% =====================================================================
\bibliographystyle{ACM-Reference-Format}
\bibliography{references}

\end{document}
\endinput
%%
%% End of file `sample-acmcp.tex'.
